(1) 质量为 $m$ 的粒子原在一维无限深阱 $0<x<L$ 的基态 $\psi_i(x)=\sqrt{2/L}\sin(\pi x/L)$。在 $t=0$ 瞬间把右壁移至 $x=2L$，其余条件不变。求新阱的本征函数与能级。
(2) 展开瞬间不变的波函数 $\psi_i(x)$ 到新阱的能量本征态，给出 $c_n$ 的通式，并特别处理 $n=2$。
(3) 求瞬间扩张后立即测量能量得到新阱第二能级的概率，以及新哈密顿量能量的平均值。
(4) 与无限缓慢绝热扩张到 $2L$ 的末态能量相比，解释两种过程为何功不同，说明概率归一与极限边界问题。
